\part{推理模型与测试时计算}

\chapter{Thinking 与 Non-Thinking 模式}
\label{sec:reasoning-models}

\begin{motivationbox}{同一组参数之外，推理阶段还能增加什么？}
传统语言模型常按固定解码策略直接生成答案；但数学、代码和规划任务需要探索中间状态、比较候选并检查约束。只增加训练参数并不是唯一提升路径：推理阶段也可以投入更多词元预算、采样次数、搜索与验证计算。推理模型关注的正是如何把这些额外计算转化为更高的问题求解质量。
\end{motivationbox}

\textbf{推理模型（reasoning model）}不是由单一网络结构定义的严格类别。面试中更稳妥的说法是：它经过数据、训练目标或推理策略优化，能够在复杂任务上有效利用较多的中间计算、搜索或验证，再输出最终答案。推理能力可能同时来自预训练中的结构学习、监督微调的示范、可验证奖励强化学习，以及推理时搜索；不能把能力归因于某一个唯一阶段。

\figReasoningModes

\section{六个容易混淆的概念}

\begin{table}[htbp]
\centering
\small
\begin{tabularx}{0.96\textwidth}{lY}
\toprule
概念 & 准确定义\\
\midrule
直接回答模式 & 提示词后立即生成面向用户的答案，通常使用较小推理预算\\
推理导向模式 & 在最终答案前投入更多内部计算、候选搜索或验证\\
思维链提示（CoT prompting） & 用提示或示例引导模型生成分步文本；它是一种交互方法，不等同于全部推理能力\\
内部推理 & 模型或系统内部用于求解的中间计算，不应假定会完整暴露给用户\\
可见推理说明 & 面向用户生成、可审核的步骤或解释；可能是内部计算的筛选、重写或单独生成结果\\
推理预算 & 对中间词元、采样次数、搜索宽度、验证调用或时间的限制\\
\bottomrule
\end{tabularx}
\caption{Reasoning 讨论中必须区分的层次。}
\end{table}

Thinking 与 Non-Thinking 可以是同一模型的两种推理行为，也可以来自不同模型。若系统只改变提示模板、最大推理词元、采样数量或验证流程，参数并未改变；若两个模式对应不同的后训练检查点、路由分支或专用模型，参数或激活路径也可能不同。因此“thinking 是否改变参数”的正确答案是：\textbf{单次模式切换不必修改参数，但产品中的两个模式不保证共享完全相同的模型实现。}

内部推理轨迹和可见解释也不能画等号。系统可以在内部进行隐藏状态计算、搜索或工具调用，只输出简洁结论；也可以额外生成一段面向用户的验证说明。可见步骤可能帮助审计，但并不保证忠实复现全部内部计算，更长的可见文本也不自动代表更强的推理。

\begin{interviewbox}{Thinking 模式和普通生成模式的本质区别是什么？}
\textbf{20 秒回答：}区别不一定在参数，而在推理阶段分配的计算与控制策略。Thinking 模式通常允许更多中间词元、多候选、搜索或验证，再产生答案；普通模式更接近单次直接生成。前者用更高延迟和成本换取复杂任务上潜在更高的正确率。

\textbf{深入解释：}语言模型的基本操作仍是条件生成。系统可以改变停止条件、推理预算、采样次数、验证器和工具闭环，使同一参数模型表现出不同推理行为；也可以通过 SFT/RL 训练它更有效地利用这些预算。模式名称本身不能说明具体训练配方。

\textbf{面试官继续追问：}
\begin{enumerate}
  \item thinking 是否意味着参数发生改变？
  \item reasoning 能力来自预训练、SFT、RL 还是搜索？
  \item CoT 与真正的问题求解能力是一回事吗？
  \item 为什么内部推理不一定对用户可见？
\end{enumerate}
\end{interviewbox}

\section{推理词元、延迟与适用任务}

设直接回答生成 $T_a$ 个词元，推理导向模式额外使用 $T_r$ 个中间词元。忽略缓存常数后，单请求解码工作随总生成长度 $T_r+T_a$ 增长；而且上下文越长，后续每步读取的键值缓存越大。若还采样 $N$ 条轨迹，粗略计算预算可写成
\[
C_{\rm test}\propto
\sum_{n=1}^{N}\bigl(T_r^{(n)}+T_a^{(n)}\bigr)
+C_{\rm verify}+C_{\rm search}.
\]
这里 $C_{\rm verify}$ 和 $C_{\rm search}$ 分别表示验证与搜索成本。真实延迟还取决于并行批处理、接受率、硬件带宽和服务调度，不能只按词元数线性外推。

数学和代码任务通常收益更明显，是因为它们有明确约束、可探索的中间结构和较可靠的验证器：答案可代入方程，代码可编译并运行测试。开放式写作没有唯一正确目标，额外搜索容易优化到狭窄评分器偏好；简单事实题的解空间又很小，长推理可能只增加延迟、引入无关分支甚至把正确直觉改错。

\begin{interviewbox}{为什么更长的 reasoning 不一定更正确？}
\textbf{20 秒回答：}更多计算只增加探索机会，不保证探索方向和验证器正确。长轨迹会累积错误、重复和无关分支；如果没有可靠反馈，模型可能用更多词元强化错误假设。收益通常先上升后饱和，甚至下降。

\textbf{深入解释：}可以把推理看成在候选解空间中分配预算。质量取决于候选覆盖、搜索策略和选择器准确率，而不是轨迹长度本身。简单问题的最优策略可能就是直接回答；复杂问题才可能值得多路径搜索和验证。系统应根据难度动态分配预算，而不是所有请求固定长思考。

\textbf{面试官继续追问：}
\begin{enumerate}
  \item reasoning token 增加会影响哪些成本？
  \item 为什么数学和代码更适合测试时扩展？
  \item 如何为简单问题设置提前停止？
  \item 长推理会怎样影响 KV Cache 和吞吐？
\end{enumerate}
\end{interviewbox}

\begin{wrongbox}
\textbf{错误说法：}能输出很长思维链，就证明模型具有更强的推理能力。

\textbf{更准确：}可见文本长度只是计算轨迹的一个表面指标。应使用最终正确率、验证通过率、鲁棒性和单位计算收益评估推理能力；可见解释还可能是对内部过程的重新组织，而非完整记录。
\end{wrongbox}

\chapter{Test-Time Scaling 与验证}
\label{sec:test-time-scaling}

\begin{motivationbox}{不增加参数，也能否用更多推理计算换质量？}
训练完成后，模型参数固定，但系统仍可增加单轨迹长度、候选数量、搜索宽度和验证次数。测试时扩展（test-time scaling）研究如何分配这些推理预算，并关注质量提升是否足以覆盖延迟和成本。
\end{motivationbox}

测试时计算的主要形式包括：

\begin{table}[htbp]
\centering
\small
\begin{tabularx}{0.96\textwidth}{lYY}
\toprule
方法 & 增加的计算 & 有效前提\\
\midrule
延长单轨迹 & 更多中间步骤 & 模型能利用长轨迹且会自我修正\\
多次采样 & 生成 $N$ 个候选 & 候选错误不完全相关\\
自洽性投票 & 对答案聚合或多数投票 & 正确答案较稳定、可规范化\\
Best-of-$N$ & 用验证器选最佳候选 & 验证器排序与真实质量一致\\
搜索 & 扩展并剪枝中间状态 & 状态可比较，分支与预算可控\\
工具验证 & 执行代码、计算或查询 & 工具反馈可靠且接口安全\\
\bottomrule
\end{tabularx}
\caption{常见测试时扩展方式及其前提。}
\end{table}

\section{多次采样、自洽性与 Best-of-N}

若单次采样正确率为 $p$，并且 $N$ 次采样近似独立，则“至少出现一个正确候选”的概率为
\[
P(\text{至少一个正确})=1-(1-p)^N.
\]
这个公式只描述\textbf{候选覆盖率}，不等于系统最终正确率。系统还必须识别哪个候选正确；若验证器不可靠，更多候选反而增加被错误高分样本欺骗的机会。实际样本也高度相关，因此独立假设常高估收益。

自洽性（self-consistency）让模型用不同采样路径得到多个答案，再对规范化后的最终答案投票。它可能有效，是因为正确推理往往汇聚到同一结果，而错误路径更分散；但若模型存在系统性偏差，多数路径会一起犯错，投票无法纠正。

Best-of-$N$ 生成 $N$ 个候选，再用奖励模型、规则或测试用例选择最高分。其前提不是“$N$ 越大越好”，而是候选分布有足够多样性、验证器与真实质量对齐、额外成本可接受。固定验证器下无限增加 $N$ 还可能产生验证器过优化：模型偶然生成专门钻评分漏洞的候选。

\begin{exercisebox}{多采样提高的是候选覆盖率}
若单次正确率 $p=0.30$，采样 $N=8$ 次，在独立假设下至少包含一个正确答案的概率为
\[
1-(1-0.30)^8=1-0.7^8\approx0.942.
\]
这不意味着最终正确率就是 $94.2\%$。若验证器只能以 $80\%$ 的概率从候选中选对，或八个样本高度相关，最终收益会明显更低。
\end{exercisebox}

\codesubsection{对应实现：Best-of-N 的系统边界}

\begin{lstlisting}
def best_of_n(prompt, model, verifier, n=8):
    # 多次 rollout 可以并行，但总生成 token 和显存占用会增加。
    candidates = [model.generate(prompt) for _ in range(n)]

    # verifier 可以是测试用例、规则或奖励模型。
    # 它必须只读取允许的信息，并返回可比较分数。
    scores = [verifier(prompt, answer) for answer in candidates]

    # 选中的是 verifier 最高分，而不必是真实质量最高的候选。
    # 因此验证器偏差会直接成为系统上限。
    best_index = max(range(n), key=lambda i: scores[i])
    return candidates[best_index], scores[best_index]
\end{lstlisting}

\begin{interviewbox}{为什么 Best-of-N 可能提高性能？}
\textbf{20 秒回答：}多次采样扩大候选覆盖，验证器再从中选择更优答案。它有效的前提是候选具有多样性，而且验证器的排序与真实质量相关；否则只是在增加重复样本或放大评分器偏差。

\textbf{深入解释：}$1-(1-p)^N$ 说明候选覆盖可随 $N$ 增长，但边际收益递减。最终性能还要乘上选择器识别正确候选的能力。验证器过弱时，增加 $N$ 可能出现“选择者诅咒”：极端高分更可能来自评分噪声或 reward hacking。

\textbf{面试官继续追问：}
\begin{enumerate}
  \item self-consistency 与 Best-of-N 的聚合方式有何区别？
  \item 样本相关性为什么会降低理论收益？
  \item verifier 不可靠会发生什么？
  \item 如何比较增加模型参数与增加测试时计算？
\end{enumerate}
\end{interviewbox}

\begin{wrongbox}
\textbf{错误说法：}测试时计算越多，质量就会单调且无限提高。

\textbf{更准确：}更多预算只提供更大的搜索机会。收益受到模型探索能力、候选相关性、验证器准确率和任务可验证性的共同限制，通常边际递减，并可能因过度思考或验证器漏洞而下降。
\end{wrongbox}

\chapter{RLVR：可验证奖励强化学习}
\label{sec:rlvr}

\begin{motivationbox}{开放式偏好难以精确，但有些任务可以自动判对错}
人类偏好奖励主观且昂贵，而数学答案、代码测试、定理检查和结构化约束常有确定性验证器。可验证奖励强化学习（RL with Verifiable Rewards，RLVR）使用这些检查结果为策略提供反馈，使模型可以通过大量在线采样探索更好的求解行为。
\end{motivationbox}

\figRLVR

对提示词 $x$，策略采样推理与答案 $y\sim\pi_\theta(\cdot\mid x)$，验证器给出
\[
r(x,y)\in\{0,1\}
\quad\text{或更细的分级分数。}
\]
数学题可验证最终数值或符号等价性，代码题可运行隔离测试，定理证明可交给形式化检查器，结构化任务可检查语法与约束。开放式写作没有唯一正确输出，很难仅用确定性检查器评价事实性、风格、创造性和帮助程度，因此通常仍需要人类偏好、奖励模型或多维规则。

\section{Outcome、Process 与稀疏奖励}

\textbf{结果监督（outcome supervision）}只检查最终答案，成本低且目标清楚，但无法区分“过程大体正确、最后算错”和“过程完全错误、偶然猜对”；模型还可能找到绕过验证器的捷径。\textbf{过程监督（process supervision）}为中间步骤提供反馈，信用分配更细，却需要可靠的步骤定义与标注，并可能把模型束缚在特定解题风格上。

若长轨迹只有最终通过时得到 $1$，绝大多数探索都得到 $0$，就形成稀疏奖励。策略很难判断哪一步值得保留，梯度估计方差也高。常见缓解思路包括：从较容易任务开始的课程学习、一次提示采样多个 rollout、提供可验证的部分分、保持足够探索熵，以及使用组内相对基线降低方差。

多 rollout 对 reasoning RL 有两层价值：一是提高偶然发现正确解的概率，二是为同一提示词提供可比较的好坏样本。GRPO 用同组奖励均值和标准差构造相对优势，正适合“一题多答”的数据形态；但若整组奖励相同，组内优势仍为零，不能凭空产生学习信号。

\begin{table}[htbp]
\centering
\small
\begin{tabularx}{0.92\textwidth}{lYY}
\toprule
监督方式 & 优点 & 风险\\
\midrule
最终结果奖励 & 自动、便宜、目标明确 & 稀疏、信用分配差、可能偶然猜对\\
过程奖励 & 步骤级反馈更密集 & 标注昂贵、步骤规范可能有偏差\\
奖励模型 & 可覆盖开放任务 & 评分偏差与 reward hacking\\
确定性验证器 & 可复现、难以主观漂移 & 只覆盖可形式化检查的维度\\
\bottomrule
\end{tabularx}
\caption{不同推理监督信号的取舍。}
\end{table}

\begin{interviewbox}{为什么数学和代码特别适合 RLVR？}
\textbf{20 秒回答：}它们常有低成本、可重复的外部验证器：数学答案可代入或做符号检查，代码可编译并跑测试。模型能大规模采样，再从明确通过/失败反馈中优化策略，而不必为每条轨迹请人打分。

\textbf{深入解释：}可靠验证器降低了奖励模型主观偏差，但最终奖励仍可能稀疏。通过测试也不一定代表实现优雅、安全或没有测试外漏洞，因此验证器覆盖面决定可学目标。开放式写作缺少唯一判据，不适合只用单一确定性奖励。

\textbf{面试官继续追问：}
\begin{enumerate}
  \item outcome reward 的信用分配问题是什么？
  \item process reward 有哪些新风险？
  \item sparse reward 为什么让探索困难？
  \item 多 rollout 与 GRPO 有什么关系？
\end{enumerate}
\end{interviewbox}

\begin{interviewbox}{验证器可靠是否就不会 reward hacking？}
\textbf{20 秒回答：}不会。验证器只检查被编码的条件，模型可能找到未覆盖的捷径，例如硬编码样例、利用测试漏洞或输出满足格式但语义错误的结果。可验证不等于验证完备。

\textbf{深入解释：}奖励优化会把策略推向验证器的边界。需要隐藏测试、对抗样例、沙箱、规则组合和人工抽检；还应监控训练奖励与独立评测是否分离。过程奖励也会被利用，只是漏洞从最终检查转移到步骤评分。

\textbf{面试官继续追问：}
\begin{enumerate}
  \item 如何设计代码任务的防泄漏测试？
  \item curriculum 为什么有助于稀疏奖励？
  \item 奖励越来越高但独立正确率下降说明什么？
\end{enumerate}
\end{interviewbox}

\chapter{推理蒸馏}
\label{sec:reasoning-distillation}

\begin{motivationbox}{昂贵推理能否转移到更小、更快的模型？}
大型教师模型可以用长轨迹、多采样和验证生成高质量训练数据，但在线部署每次都承担同样预算过于昂贵。推理蒸馏尝试把教师发现的解题模式压缩到学生参数中，使学生用更少测试时计算复现部分能力。
\end{motivationbox}

\figReasoningDistill

普通 logits 蒸馏让学生匹配教师的软类别分布；回答蒸馏用教师生成的输入--输出对训练学生；推理蒸馏还可能使用经过过滤的中间步骤或结构化求解轨迹。只蒸馏最终答案数据简单、暴露错误过程较少，但学生看不到如何到达答案；蒸馏轨迹能提供更密集监督，却会把教师的冗余、风格偏差和错误推理一起传递。

推理蒸馏能让小模型变快，是因为部分搜索工作被“摊销”进训练：教师离线探索许多路径，学生在参数中学习常见映射，在线时不必重复同样规模的搜索。但学生容量、训练分布和验证过滤决定上限；蒸馏不能保证小模型在分布外任务上复现教师的全部探索能力。

\begin{interviewbox}{为什么 reasoning 能力可以蒸馏给小模型？}
\textbf{20 秒回答：}教师用高预算产生并验证求解示范，学生通过监督学习把这些常见推理模式压缩进参数。在线时学生可用更短轨迹直接复现部分解法，因此延迟更低；但它继承教师和过滤器的偏差，也受学生容量限制。

\textbf{深入解释：}最终答案蒸馏提供硬目标，轨迹蒸馏提供步骤级条件分布，logits 蒸馏还保留候选间相对信息。高质量过滤十分关键：错误但流畅的轨迹会成为强监督噪声，甚至让学生学会合理化错误答案。

\textbf{面试官继续追问：}
\begin{enumerate}
  \item 推理蒸馏与普通 logits 蒸馏有什么区别？
  \item 只蒸馏答案和蒸馏轨迹各有什么取舍？
  \item 为什么蒸馏模型可能更短、更快？
  \item teacher 的错误 reasoning 会发生什么？
\end{enumerate}
\end{interviewbox}

\begin{pitfallbox}
不要把“教师生成了长解释”直接等同于“高质量推理数据”。商业训练流水线还需要答案验证、轨迹去重、难度配比、污染检测和安全过滤；具体私有模型的训练配方不能从公开输出反推。
\end{pitfallbox}
